% Cuadrícula de 5 mm con márgenes (A5) — Eric Lucero González — https://ericlucero.dev
% Licencia CC BY-NC-SA 4.0. Compilar con pdfLaTeX (dos pasadas: latexmk lo hace solo).
\documentclass[a5paper]{article}
\usepackage[T1]{fontenc}
\usepackage[spanish]{babel}
\usepackage{lmodern}
\usepackage{tikz}
\pagestyle{empty}

% ── Configuración (mm) ─────────────────────────────────────
\def\anchoPapel{148}
\def\altoPapel{210}
\def\margenLateral{4}   % margen mínimo a izquierda y derecha
\def\abajo{12}          % franja libre abajo
\def\arriba{13}         % franja libre arriba
\def\paso{5}           % lado de cada cuadro
\def\cuadrosIzq{3}      % la línea de margen izquierda, a tantos cuadros del borde de la rejilla
\def\cuadrosDer{3}      % ídem, línea derecha

\begin{document}
\begin{tikzpicture}[
    remember picture, overlay, shift={(current page.south west)},
    x=1mm, y=1mm,
    grid lines/.style={gray!45, very thin},
    margin lines/.style={red!50!gray, thin}
  ]
  % Cuadros enteros que caben; centrada en horizontal, empieza en \abajo
  \pgfmathtruncatemacro{\nx}{(\anchoPapel - 2*\margenLateral) / \paso}
  \pgfmathtruncatemacro{\ny}{(\altoPapel - \abajo - \arriba) / \paso}
  \pgfmathsetmacro{\ancho}{\nx * \paso}
  \pgfmathsetmacro{\alto}{\ny * \paso}
  \pgfmathsetmacro{\x}{(\anchoPapel - \ancho) / 2}

  \draw[step=\paso, grid lines, shift={(\x,\abajo)}] (0,0) grid (\ancho,\alto);

  % Líneas de margen, sobre líneas de la rejilla y a toda altura
  \pgfmathsetmacro{\xIzq}{\x + \cuadrosIzq*\paso}
  \pgfmathsetmacro{\xDer}{\x + \ancho - \cuadrosDer*\paso}
  \draw[margin lines] (\xIzq, 0) -- (\xIzq, \altoPapel);
  \draw[margin lines] (\xDer, 0) -- (\xDer, \altoPapel);
\end{tikzpicture}
\end{document}
